\section{Novelty is bounded and defeasible}

A verified theorem is not automatically new. Newness requires a literature world, an equivalence notion and a cutoff. Define a bounded novelty certificate
\[
N_t(T)=(C_{\le t},S,\nu,f(T),R),
\]
where $C_{\le t}$ is a named corpus snapshot, $S$ the search routes, $\nu$ a normalization/equivalence procedure, $f(T)$ a theorem fingerprint and $R$ the review record. Search routes should include exact wording, terminology variants, notation normalization, structural similarity, stronger-parent theorem retrieval, citation neighborhoods and multilingual/domain-synonym searches.

\begin{proposition}[Novelty non-monotonicity]
Let $C_t\subseteq C_{t+1}$ be literature corpora and define $\mathrm{Novel}_{C}(T)=1$ when no equivalent prior result is found in $C$. Then
\[
\mathrm{Novel}_{C_t}(T)=1
\]
does not imply
\[
\mathrm{Novel}_{C_{t+1}}(T)=1.
\]
\end{proposition}

\begin{proof}
Choose a theorem $T'$ equivalent to $T$ such that $T'\notin C_t$ and $T'\in C_{t+1}$. The novelty predicate is true at $t$ and false at $t+1$. \qed
\end{proof}

Thus novelty authority is defeasible even when proof validity of the fixed statement is unchanged. The licensed statement is ``no equivalent result was found in the registered novelty world at cutoff $t$,'' not ``global novelty has been proved.'' This distinction also explains why rediscovery can be scientifically informative without being mislabelled as new mathematics.
